%% ma: L12-B08-0094
%% lop: 12
%% chuong: 2
%% bai: 8
%% dang: 12.8.29
%% loai: TN4
%% muc: NB
%% dap_an: "B"
%% nguon: "CĐ2.3.pdf giáo viên gửi 27/10/2026 (M07, Lê Minh Khôi – Nguyễn Anh Khôi, Chương II Chủ đề 3 Biểu thức tọa độ của các phép toán vectơ), Câu 2"
%% kien_thuc: "Bài 8 (biểu thức tọa độ của phép cộng, trừ vectơ, tích vectơ với một số)"
%% trang_thai: chinh-thuc
\begin{ex}
Trong không gian $Oxyz$, cho vectơ $\vec a=(-2;6;2)$. Tọa độ của vectơ $\vec m=\dfrac32\cdot\vec a$ là
\choice{$\vec m=\left(-\dfrac43;4;\dfrac43\right)$}{\True $\vec m=(-3;9;3)$}{$\vec m=(6;9;6)$}{$\vec m=(-3;6;3)$}
\loigiai{$\vec m=\left(\dfrac32\cdot(-2);\dfrac32\cdot6;\dfrac32\cdot2\right)=(-3;9;3)$.}
\end{ex}
